\clearpage
\section{Allocation successor coherence and exact costs}
\label{sec:copying-allocation-coherence}
The allocation branch of matching independence is now complete for the actual
combined rewrite relation. If two rewrites start at the same ref request in
the same host graph, their successor graphs are recursively isomorphic and
the isomorphism carries one resumed token to the other. The matches may choose
different retained hosts, stable argument subgraphs, context decompositions
and output representations. No agreement of these choices is assumed.
The theorem covers every finite interface and depth allowance, including
arbitrary typed context insertion and simultaneous expansion.

This is stronger than source selection: the entire result is compared,
including the fresh instance, canonical atom tree, store, retained graph and
token. It completes the allocation case of the copying-machine determinism
obligation. Section~\ref{sec:copying-equality-coherence} subsequently completes
both reference-equality branches. The remaining successor families,
refocusing and the concrete specimen interface remain open. Neither the
Proposition~12.2 terms nor Figure~16's translation is modified.

The construction follows the allocation rule in PDF Figure~19(a). Its source
is a ref corolla followed by a stable argument $A$; its target replaces that
corolla by the allocated-cell prefix and keeps $A$. The proof turns this
observation into a canonical replacement of one edge of the entire source
host. It uses the previously proved substitution associativity and unit laws,
then inducts over the actual context-closed rewrite derivation. The original
stability premise is retained: a ref head alone need not be reducible.

The dependency chain is therefore source-head uniqueness, state-level
replacement congruences, canonical local allocation, induction over actual
matches, successor coherence, then exact successful-cost consequences.
This discharges a concrete machine obligation needed before the paper's
robustness theorem can be instantiated; it is not yet the reference-creation
template's robustness proof.

\Needspace{12\baselineskip}
\subsection{States and replacement congruences}

A state consists of a graph, a focus mode and a vertex. When a graph is
changed by equality, transporting its vertex along that equality gives the
same dependent state. Equality of states supplies a state isomorphism by the
identity construction. Both facts follow by eliminating the equality proof;
they handle representation indices without dropping the token coordinate.

\lean{Copying/ReplacementStates}{12}{15}{Combined.StateIso.cast_state}

\lean{Copying/ReplacementStates}{28}{29}{Combined.StateIso.ofEq}

A state isomorphism between two fillers lifts through plugging into a fixed
host edge. The existing graph congruence gives the graph isomorphism; its
inserted-port equation, followed by the original token equation, gives the
new token equation. For simultaneous expansion, assign the supplied graph
isomorphism to the changed component and identity isomorphisms to all other
components. The assembly isomorphism preserves the selected token by its
component-insertion equation. These statements do not assume that attachment
is injective on boundary ports.

\lean{Copying/ReplacementStates}{18}{25}{Combined.StateIso.plug}

\lean{Copying/ReplacementStates}{32}{47}{Combined.StateIso.expand}

Changing equal interface lists preserves a primitive-head witness. Labels
are unchanged and the source incidence list is mapped by the boundary cast;
source membership therefore survives by membership in a mapped list.

\lean{Copying/ReplacementStates}{51}{54}{Combined.head_reboundary}

\termdef{canonical-replacement-state}
For a host $G$, selected edge $e$, fitting replacement $A$ and old token $v$,
define the \emph{replacement state}
\[
 R(G,e,A,v)=\bigl(G[e:=A],\;\mathrm{search},\;\iota_G(v)\bigr),
\]
where $\iota_G$ is the actual retained-host port map. The old wire is retained,
including when it is a boundary wire or participates in cyclic host structure.
This definition records where search resumes; it does not itself assert that
the operational relation can perform the replacement.

\lean{Copying/ReplacementStates}{58}{61}{Combined.Replacement.state}

Replacing an edge of a plugged filler can be performed either after
plugging or inside the filler before plugging. The graph isomorphism is
substitution associativity; its middle-component equation carries the old
filler token to the same attached wire. In a simultaneous expansion, regroup
the replacement into its selected component and leave every other component
unchanged. The expansion/replacement interchange isomorphism provides the
graph map and its component equation provides the token map. The cast in the
formal statement only identifies the selected component of the updated
family with the explicitly replaced filler.

\lean{Copying/ReplacementStates}{64}{73}{Combined.Replacement.plug}

\lean{Copying/ReplacementStates}{76}{88}{Combined.Replacement.expand}

Recursive renaming of the host transports replacement to the renamed edge.
The replacement's interface is reindexed by the equalities forced by the
related edge labels. Host-plug congruence gives the graph isomorphism, and
its retained-host equation proves the token equation. Finally, if a changed
component has an independently chosen isomorphic output, compose expansion
interchange with expansion congruence. Equality transport removes the selected
family-component cast before applying that component isomorphism.

\lean{Copying/ReplacementStates}{91}{101}{Combined.Replacement.rename}

\lean{Copying/ReplacementStates}{104}{121}{Combined.Replacement.expand_result}

\Needspace{12\baselineskip}
\subsection{Canonical allocation in a source host}

Two successive changes of interface lists equal the single change by their
composite equalities. Eliminating the four equality proofs reduces the claim
to the identity-boundary law. In an allocation source, the selected head is
the left serial component's sole ref edge. Its primitive label and source
membership follow directly from the ref corolla and serial incidence.

\lean{Copying/AllocationCoherence}{9}{12}{reboundary_comp}

\lean{Copying/AllocationCoherence}{16}{17}{Allocation.head}

\lean{Copying/AllocationCoherence}{19}{21}{Allocation.head_at}

\termdef{canonical-allocation-successor}
Suppose $e$ has label ref in $G$. Its interfaces are then both the singleton
term interface. The \emph{allocation filler} is the existing allocated-cell
prefix, reindexed by these forced interface equalities. It contains an
instance edge, the original canonical atom-sharing tree and a store edge.
The \emph{canonical allocation successor} is
\[
 C(G,e,v)=R(G,e,\mathrm{cellPrefix},v).
\]
Thus the construction depends only on the source host, its selected edge and
its focused wire. It does not inspect a chosen decomposition into argument
and retained context, nor simplify any distributor.

\lean{Copying/AllocationCoherence}{24}{29}{Allocation.filler}

\lean{Copying/AllocationCoherence}{32}{34}{Allocation.canonical}

Canonical allocation commutes with recursive host renaming. Apply the
replacement-renaming isomorphism and combine the successive interface
reindexings; proof irrelevance identifies the resulting equality witnesses.
For a local source $\mathrm{ref};A$, serial replacement interchange first moves
the replacement inside the left ref corolla. The corolla substitution unit
then identifies that replaced component with the allocated-cell prefix.
Serial congruence leaves $A$ fixed. All these isomorphisms fix the external
input, so the token equation follows as well.

\lean{Copying/AllocationCoherence}{37}{48}{Allocation.canonical_rename}

\lean{Copying/AllocationCoherence}{51}{62}{Allocation.canonical_local}

\Needspace{12\baselineskip}
\subsection{Independently chosen contextual matches}

Every actual rewrite already has a recognized primitive head in its changed
component. Map that head into a surrounding occurrence. If the host token
has primitive head $e$, source uniqueness identifies the mapped component
head with $e$, and primitive-label injectivity identifies their labels.
Consequently the component contains the same primitive head. This reflection
argument uses the actual rewrite witness; it does not claim that arbitrary
occurrences reflect all source information, and needs no injective port map.

\lean{Copying/AllocationCoherence}{71}{76}{Combined.rewrite_head_reflect}

For a local rule whose source head is ref, inspect the actual local-rule
constructor. Allocation has the canonical result just proved. Dereference,
assignment, both equality branches, application, term copying and arithmetic
have different primitive heads, contradicting head uniqueness. The allocation
constructor still carries its original stable-argument proof. Thus the local
case contains neither an assumed determinism law nor an extra reduction rule.

\lean{Copying/AllocationCoherence}{79}{113}{Combined.local_allocation_canonical}

Now induct over the actual contextual rewrite derivation. The local case is
the preceding result. In a renaming case, reflect the ref head to the original
source, apply induction, and compose the inverse canonical-source renaming
with the inductive result and target renaming. In a plugging case, reflect
the head into the changed filler, apply induction, and compose replacement
associativity with state congruence. In an expansion case, reflect into the
selected component and use the proved expansion-result interchange. In each
case the head's edge equality identifies the canonical replacement of the
whole source with the one used in the induction.

\lean{Copying/AllocationCoherence}{117}{138}{Combined.rewrite_allocation_canonical}

\termdef{allocation-successor-coherence}
\emph{Successor coherence at allocation} means that any two actual rewrites
from the same ref request produce state-isomorphic successors. Both successors
are isomorphic to $C(G,e,v)$; invert the first isomorphism and compose it with
the second. This compares complete output graphs and resumed token positions,
not merely labels or observable answers.

\lean{Copying/AllocationCoherence}{141}{147}{Combined.allocation_rewrite_deterministic}

The corresponding full-machine step cannot be a search step: search has no
transition from rewrite mode. Therefore every full step at the request has
the canonical successor, and any two such steps have isomorphic successors.
Moreover, an admitted allocation can step directly to the canonical
representative: rename the target of its actual rewrite by the inverse
canonical-result isomorphism. Its token equation gives the required resumed
coordinate. The existence premise here is an actual transition, not a
hypothetical stability or correctness axiom.

\lean{Copying/AllocationCoherence}{155}{160}{Combined.allocation_step_canonical}

\lean{Copying/AllocationCoherence}{163}{169}{Combined.allocation_step_deterministic}

\lean{Copying/AllocationCoherence}{172}{184}{Combined.allocation_step_to_canonical}

\Needspace{12\baselineskip}
\subsection{Exact costs and the unbounded machine}

Write $T_n(s)$ for existence of a successful run of exactly $n$ combined
transitions from $s$. Given an actual allocation step $s\to t$,
\[
 T_{n+1}(s)\quad\Longleftrightarrow\quad T_n(t).
\]
For the forward direction, split off the successful run's first step. Its
successor is state-isomorphic to $t$, and the previously proved run transport
carries the tail at the same cost. Conversely, prepend the given allocation
step. This argument places no restriction on the later computation; in
particular it does not assume termination of a callback.

\lean{Copying/AllocationCoherence}{188}{199}{Combined.allocation_cost_iff}

An unbounded step retains its finite representation allowance. Inverting its
wrapper reduces two allocation steps to the bounded result above, then
re-wrapping the state isomorphism gives the existing representation relation.
That relation preserves every successful cost. Thus independently chosen
unbounded allocation successors refine one another for every cost preorder,
including exact equality. This remains an allocation theorem, not global
unbounded-machine determinism.

\lean{Copying/AllocationCoherence}{202}{211}{Combined.allocation_unbounded_deterministic}

\lean{Copying/AllocationCoherence}{214}{218}{Combined.allocation_unbounded_equiv}

An actual search step has a unique full-machine successor: another search
step agrees by intrinsic search determinism, while a rewrite alternative
would require rewrite mode, from which the given search step cannot start.
A final state has precisely cost zero, because finality excludes any first
transition. Its zero-step run supplies the converse.

\lean{Copying/AllocationCost}{8}{12}{Combined.search_successor_unique}

\lean{Copying/AllocationCost}{15}{20}{Combined.final_cost_iff}

More generally, a given step $s\to t$ whose alternative successors are all
state-isomorphic to $t$ satisfies
\[
 T_n(s)\quad\Longleftrightarrow\quad
 \exists m\in\mathbb N\;\bigl(n=m+1\ \land\ T_m(t)\bigr).
\]
A successful zero-step run at $s$ would contradict its outgoing step. For a
positive run, remove its first step and transport its tail by successor
coherence. The reverse direction prepends the given step. Search uniqueness
and proved allocation coherence discharge this intermediate theorem's local
coherence hypothesis in their respective specializations. The public search
and allocation cost laws have no assumed determinism premise.

\lean{Copying/AllocationCost}{23}{34}{Combined.coherent_step_cost_iff}

\lean{Copying/AllocationCost}{37}{40}{Combined.search_cost_iff}

\lean{Copying/AllocationCost}{43}{47}{Combined.allocation_cost_spectrum}

\Needspace{12\baselineskip}
\subsection{Regressions on the original translation and arbitrary hosts}

The literal translated term $\mathtt{ref}\ n$ retains all Figure~16
single-use distributors, including its empty-environment forests. Its ref
head has the required label and input incidence by definition. The actual
allocation transition is obtained by lifting the existing stable-argument
allocation through the unchanged serial distributor continuation. Applying
successor coherence compares every other full-machine successor with this
constructed result.

\lean{Copying/AllocationCoherenceExamples}{9}{10}{AllocationCoherenceExamples.reference_head}

\lean{Copying/AllocationCoherenceExamples}{13}{17}{AllocationCoherenceExamples.reference_allocate}

\lean{Copying/AllocationCoherenceExamples}{20}{23}{AllocationCoherenceExamples.reference_successor}

For every natural $n$ and cost $k$, the original translation of
$\mathtt{ref}\ n$ satisfies
\[
 T_k(\lbrack\!\lbrack\mathtt{ref}\ n\rbrack\!\rbrack,\mathrm{search},\mathrm{input}_0)
 \quad\Longleftrightarrow\quad k=5.
\]
There are three search transitions (enter ref, return the numeral, request
allocation), one allocation transition and one instance-return transition.
Apply the exact cost laws successively, then finality's zero-cost law. The
resulting five successor equations force $k=5$; the reverse direction supplies
those five costs. This proves uniqueness of successful cost under the full
combined relation, beyond the earlier exhibited five-step run.

\lean{Copying/AllocationCoherenceExamples}{27}{42}{AllocationCoherenceExamples.reference_exact_cost}

Exact bounded/unbounded run correspondence gives the same cost characterization
for the unbounded initializer. The already proved padding reflection and
preservation then give the same characterization at every larger allowance.
In particular this covers the literal $\mathtt{ref}\ 1$ subexpression used
in Proposition~12.2. It does not claim the proposition's entire closure is
observationally equivalent yet.

\lean{Copying/AllocationCoherenceExamples}{45}{54}{AllocationCoherenceExamples.reference_unbounded_exact_cost}

\lean{Copying/AllocationCoherenceExamples}{57}{60}{AllocationCoherenceExamples.reference_padded_exact_cost}

Finally, insert an allocation source with any stable argument into any
typed host edge whose input interface is a single term. Reindex its boundary
by the supplied interface equality and lift the actual local allocation step.
Map the ref-head witness through the same attachment. The general coherence
theorem shows that every alternative successor agrees with the host filled
by the allocation result. No acyclicity, global boundary injectivity or
common matching decomposition is required.

\lean{Copying/AllocationCoherenceExamples}{64}{74}{AllocationCoherenceExamples.host_allocation_successor}

All 39 new declarations are audited. The full build passes 373 jobs and the
three audits inspect 6 abstract, 477 graph and 1143 copying declarations,
with only \texttt{propext}, \texttt{Classical.choice} and \texttt{Quot.sound}.
The standard check script now includes the copying audit as well.

Appendix~B.1 removes allocation from the unfinished successor-coherence
families. Its four broad packages remain open: in particular copying needs
comparison of complete payloads, residual trees, distributors and retained
hosts, and all eleven template families still need the required robustness
arguments before Figure~22 can be assembled.

The earlier estimate of 80\% completion is withdrawn. A cautious current
assessment is \textbf{50--65\% of the mathematical effort}, with low confidence
(26 October 2026). That is not a measured percentage, a count of declarations,
or a claim that the remaining work is routine. Source selection and allocation
coherence are substantial established prerequisites; full copying-machine laws,
the dependent template proofs and the final observational argument can still
account for a large part of the work. This milestone does not justify a
numerically precise increase within that range.

